%%FILE XIIc-a
%%BLOCK after="Intermediate categories and the extra information in a replacement"
\begin{dcomment}{development}{p218}{Axioms (a) and (b) do not imply (c): a lattice with six elements}
The question of ms.~217--218 has a negative answer. The three assertions are read as follows, for a category $\mathcal M$ with pullbacks, a class $W$, the full subcategory $\mathcal F_S\subset\mathcal M/S$ of the intrinsic $W$-fibrations over $S$, and the class $W_S$ of the arrows over $S$ that are in $W$: (a) the inclusion of $\mathcal F_S$, localized at the arrows that remain in $W$ after every base change, into $\mathcal M/S$ localized at $W_S$ is an equivalence; (b) an arrow of $W_S$ between objects of $\mathcal F_S$ remains in $W$ after every base change; (c) for every full subcategory $\mathcal H$ with $\mathcal F_S\subset\mathcal H\subset\mathcal M/S$ the two inclusions induce equivalences of the localizations at the restrictions of $W_S$. The example was not found in the sources consulted; no priority is claimed.

\textbf{Proposition.}\ours
\textit{Let $\mathcal M$ be the ordered set $\{0,\dots,5\}$ generated by $0<1$, $0<2$, $1<3$, $2<3$, $1<4$, $3<5$, $4<5$, and let $W$ consist of the identities and of the arrows $1\to3$, $1\to4$, $2\to3$. Then $\mathcal M$ is a lattice, $W$ is strongly saturated, (a) and (b) hold over every base, and (c) fails over the base $5$ for $\mathcal H=\{0,3,4,5\}$.}

\textit{Proof.}
Pullbacks are meets: the meet of two comparable elements is the smaller one, and $1\wedge2=2\wedge4=0$, $3\wedge4=1$. An arrow $x\leq s$ is an intrinsic fibration if and only if $a\wedge x\to b\wedge x$ is in $W$ for every $a\to b$ in $W$ with $b\leq s$. This gives $\mathcal F_s=\mathcal M/s$ for $s=0,1,2,4$, $\mathcal F_3=\{0,3\}$ and $\mathcal F_5=\{0,3,5\}$: over $3$ the element $1$ fails on $2\to3$ (which becomes $0\to1$) and $2$ fails on $1\to3$; over $5$ moreover $4$ fails on $2\to3$; and for $x=3$ the three weak arrows become $1\to3$, $\mathrm{id}_1$, $2\to3$.

(b). For $s\neq4$ the arrows of $W$ between objects of $\mathcal F_s$ are identities. For $s=4$ there is also $1\to4$, whose pullbacks along $0,1,4$ are $\mathrm{id}_0$, $\mathrm{id}_1$ and itself.

(a). Only the bases $3$ and $5$ need an argument. Over $3$ the map $R_3$ with $R_3(0)=0$ and $R_3(1)=R_3(2)=R_3(3)=3$ is a retraction of $\mathcal M/3$ onto $\mathcal F_3$ that preserves $W$, with $x\leq R_3(x)$ in $W$ for all $x$. Over $5$ let $R$ and $T$ be the monotone maps
\[
\begin{array}{c|cccccc}
x&0&1&2&3&4&5\\ \hline
R(x)&0&3&3&3&3&5\\
T(x)&0&1&2&3&1&5 .
\end{array}
\]
Both preserve $W$; $T(x)\leq x$ and $T(x)\leq R(x)$ for all $x$, with all these arrows in $W$; and $R$ is the identity on $\mathcal F_5$. So after localization the inclusion of $\mathcal F_5$ and $R$ are inverse equivalences.

Strong saturation. An arrow $a\to b$ is in $W$ exactly when $R(a)=R(b)$. So $R$ factors through the localization, and an arrow that becomes invertible there is sent by $R$ to an identity of the chain $\{0<3<5\}$.

(c) fails. The restrictions of $W_5$ to $\mathcal H=\{0,3,4,5\}$ and to $\mathcal F_5$ contain only identities, so the first localized inclusion is the inclusion of the chain $0<3<5$ in the ordered set $\mathcal H$, which is not essentially surjective. The second is not full: there is no arrow from $3$ to $4$ in $\mathcal H$, and $3\leftarrow1\to4$ is an isomorphism in the localization of $\mathcal M$.\hfill$\square$

The zigzag $x\leftarrow T(x)\to R(x)$ passes through $T(4)=1$, which is not in $\mathcal H$: this is why it gives (a) and not (c). In the example the arrow $4\to5$ has no factorization as an arrow of $W$ followed by an intrinsic fibration. With a functorial replacement $R_S:\mathcal M/S\to\mathcal F_S$ and a natural transformation from the identity with components in $W_S$, (c) holds, as shown in the preceding commentary. Whether (c) follows from (a), (b) and the existence of factorizations, without functoriality, is not decided here.
\end{dcomment}

%%APPEND title="The saturation question has an affirmative answer in Cat"
The equality $\widehat{\mathcal L}=\mathrm{Liss}_W$ holds for every fundamental localizer satisfying the relative axiom, with no accessibility hypothesis, and dually the functors that universally preserve the class defined by the tests $s\backslash X$ are the $W$-proper ones. The argument for the inclusion $\mathrm{Liss}_W\subset\widehat{\mathcal L}$ uses, for $X$ over $T$, the opfibration $R_TX=\int_{t}X/t\to T$ and the functor $j_X:X\to R_TX$, $x\mapsto(p(x),x,\mathrm{id})$, whose comma categories are the $X/x$, so that $j_X$ is aspherical; pullback along a smooth functor preserves aspherical functors and fibrewise equivalences of opfibrations, and two-out-of-three in the square formed by $j_X$, $j_Y$ and their pullbacks concludes.

%%BLOCK after="Smoothness and properness depend on a coefficient theory"
\begin{dcomment}{development}{p225}{Complementarity holds for ordinary weak equivalences only}
The inclusion $\mathrm{Liss}\cap\mathrm{Prop}\subseteq\mathrm{Fib}_W$ considered on this page, with $\mathrm{Liss}$ and $\mathrm{Prop}$ the $W$-smooth and $W$-proper functors, characterizes the two extreme localizers. The statement was not found in this form in the sources consulted; no priority is claimed.

\textbf{Theorem.}\ours
\textit{Let $W$ be a fundamental localizer satisfying the relative axiom. Every functor that is $W$-smooth and $W$-proper is an intrinsic $W$-fibration if and only if $W=W_\infty$ or $W$ is the class of all functors.}

\textit{Proof.}
For $W_\infty$ a strict pullback along a smooth and proper functor is a homotopy pullback \citep[6.4.8 and 6.4.16]{Cisinski2006}, so the functor preserves weak equivalences under every base change. For the class of all functors there is nothing to prove.

Conversely suppose the inclusion holds for $W$. Since $W_\infty\subset W$, a $W_\infty$-smooth functor is $W$-smooth and a $W_\infty$-proper functor is $W$-proper. Every functor $v:C\to B$ factors as a weak equivalence $C\to P$ followed by a functor $p:P\to B$ that is $W_\infty$-smooth and $W_\infty$-proper (the path factorization; commentary on Chapter~VII, ms.~157). By hypothesis $p$ is in $\mathrm{Fib}_W$, and strict pullbacks along $p$ are homotopy pullbacks. So for $u:A\to B$ in $W$ the functor $A\times_BP\to P$ is in $W$, and it represents the homotopy pullback of $u$ along $v$: the class $W$ is stable under homotopy base change. Let $u\in W$ be not in $W_\infty$. Some homotopy fibre $F$ of $u$ is not weakly contractible and is $W$-aspherical. If $F$ is empty, $\emptyset\to e$ is in $W$. Otherwise every point $e\to F$ is in $W$, and pulling a point back along itself shows that the loop space of $F$ is $W$-aspherical; repeating, one reaches a $W$-aspherical space with at least two components, and the homotopy pullback of two points in different components is empty. So $\emptyset\to e$ is in $W$ in all cases; then $\emptyset\to A$ is in $W$ for every $A$, by stability under products, and every functor is in $W$.\hfill$\square$

Examples for the localizers between the two: the functor from the contractible groupoid with two objects to the group with two elements for $W_0$, and the path functor over a finite category of the homotopy type of $S^{n+1}$ for the $n$-equivalences, as in the preceding commentary. A functor that is a fibration and an opfibration serves for $W_1$: the central extension $E_c\to B$ by $C_2$ of the ordered set $B$ of faces of the boundary of the tetrahedron, twisted by a cocycle supported on one triangle (commentary on Chapter~I, ms.~122). It is smooth and proper for every $W$. The functor $e\to B$ is a $W_1$-equivalence, because $B$ has the type of $S^2$; its pullback is the inclusion of the fibre $BC_2$ in $E_c$, which is simply connected, so the pullback is not a $W_1$-equivalence. The same functor is stated to serve for the localizer of locally constant rational coefficients; that verification is not reproduced here. For $W_\infty$ there are strict inclusions: fibrations that are opfibrations, then smooth and proper functors, then intrinsic fibrations; the first is strict by the four-element ordered set of the commentary on Chapter~VI, ms.~186.
\end{dcomment}

%%FILE XIIc-b
%%APPEND title="The factorization axiom and Brown's hypothesis"
What the left factorization axiom gives can be stated exactly. With a terminal object and pullbacks, if every arrow is an arrow of $W$ followed by an intrinsic $W$-fibration, then the arrows of $W\cap\mathrm{Fib}_W$ remain in $W$ after every pullback, the objects that are fibrant over the terminal object form a category of fibrant objects in the sense of Brown which presents the same localization, and finite homotopy limits and the replacement of finite inverse diagrams by fibrant ones follow. It does not give all small homotopy Kan extensions: finite sets with the isomorphisms satisfy the finite axioms. Nor does it give the two fine theories over a base, which need a choice of test maps.

%%FILE XIII-b
%%BLOCK after="Why uniform left structures produce right images"
\begin{dcomment}{development}{p57}{A closed exact structure is a closed model structure}
The uniform hypothesis of this page asks for structures that are closed and exact. Together the two adjectives give more than either: the classes of trivial cofibrations and trivial fibrations are then the intersections with $W$, as in a model category in the sense of Quillen with finite limits and colimits.

\textbf{Proposition.}\ours
\textit{Let $(C,TC,F,TF)$ satisfy the axioms of ms.~27--28: $W$ has two-out-of-three, every arrow factors as $TC$ followed by $F$ and as $C$ followed by $TF$, and $TC\subset C\cap W$, $TF\subset F\cap W$. Suppose that the two lifting pairs $(C,TF)$ and $(TC,F)$ are closed and that $W=TF\circ TC$. Then $TC=C\cap W$, $TF=F\cap W$, and $W$ is stable under retracts.}

\textit{Proof.}
Let $i:A\to B$ be in $C\cap W$, and write $i=pj$ with $j$ in $TC$ and $p:Z\to B$ in $TF$. Since $i$ lifts against $p$, there is $s:B\to Z$ with $si=j$ and $ps=\mathrm{id}_B$; so $i$ is a retract of $j$, and $TC$, defined by a lifting property, is stable under retracts. Dually $TF=F\cap W$.

Let $f:A\to B$ be a retract of $w:X\to Y$ in $W$, with sections $a,b$ and retractions $r,s$. Suppose first $f\in F$. Factor $w=vj$ with $j\in TC$ and $v\in TF$; a lift of $j$ against $f$ in the square given by $r$ and $sv$ exhibits $f$ as a retract of $v$, so $f\in TF$. In general factor $f=hg$ with $g:A\to Z$ in $TC$ and $h$ in $F$, and form the pushout $P$ of $X\leftarrow A\to Z$ (maps $a$, $g$), with $i:X\to P$ in $TC$ and $x:Z\to P$. The maps $w$ and $bh$ induce $j:P\to Y$ with $ji=w$, so $j\in W$; the maps $gr$ and $\mathrm{id}_Z$ induce $y:P\to Z$ with $yx=\mathrm{id}_Z$ and $hy=sj$. Thus $h$ is a retract of $j$, hence in $W$ by the first case, and $f=hg$ is in $W$.\hfill$\square$

So under the uniform hypothesis each diagram category carries a closed model structure in the finite sense, proper when the conditions M$''$ hold, and the argument of Brown in the next commentaries applies to every functor whose ordinary right Kan extension exists. An exact structure that is not closed, or a structure exact on one side only, does not have these consequences.
\end{dcomment}

%%APPEND title="The proposed noetherian enlargement is actually finite"
The same conclusion has been reached independently: an ordered set whose Alexandrov topology is noetherian and sober is finite, by an argument with closed sets and Ramsey's theorem.

%%FILE XVI
%%BLOCK after="Axioms for equivalences before a choice of model"
\begin{dcomment}{development}{p2}{The absolute axiom does not imply the relative one}
Definition~1.1 asks for Loc~(1), (2), (3), the last being the absolute form of Theorem~A; the relative form Loc~(3 bis) is introduced on ms.~3, where Loc~(3) is noted to be its special case. The converse fails for classes satisfying only the moderate saturation of Loc~(1). The following example has been checked in full; whether it, or an equivalent one, is in the literature on weak fundamental localizers has not been checked, and no claim of novelty is made.

\textbf{Proposition.}\ours
\textit{Let $W^\dagger$ be the class of functors $f:A\to B$ such that $\pi_0(f)$ is bijective and, if $A$ is connected, $f$ is a weak equivalence. Then $W^\dagger$ satisfies Loc~(1), (2), (3) and does not satisfy Loc~(3 bis). It is not stable under retracts.}

\textit{Proof.}
If two of $f$, $g$, $gf$ are in $W^\dagger$ the three functors are bijective on components and the three categories have the same number of components; when that number is one, two-out-of-three for $W_\infty$ applies. If $gf$ and $fg$ are in $W^\dagger$, then $\pi_0(f)$ and $\pi_0(g)$ are bijective, and in the connected case $f$ and $g$ are homotopy equivalences of classifying spaces because both composites are. This is Loc~(1). A category is $W^\dagger$-aspherical if and only if it is weakly contractible; so Loc~(2) holds, and if every $A/b$ is $W^\dagger$-aspherical then $f$ is in $W_\infty\subset W^\dagger$ by Quillen's Theorem~A, which is Loc~(3).

Let $G$ be the group with two elements, $E$ the groupoid with the two objects $x\in G$ and exactly one arrow between any two of them, and $p:E\to BG$ the functor sending the arrow $x\to x+g$ to $g$. Put $A=E\times BG$, $B=E$, $u:A\to B$ the projection, both regarded over $BG$ through $p$. The category $E$ is contractible and $A$ is connected with automorphism groups $G$, so $u$ is not a weak equivalence, and not in $W^\dagger$. Over the unique object $c$ of $BG$, the category $E/c$ has the four objects $(x,t)$ and an arrow $(x,t)\to(x+g,t')$ exactly when $t'+g=t$; so $t+x$ is constant along arrows and $E/c$ has two components, each contractible. Moreover $A/c=(E/c)\times BG$ and $u/c$ is the projection, a bijection on the two components. Hence $u/c\in W^\dagger$ and $u\notin W^\dagger$: Loc~(3 bis) fails.

The functor $BG\to e$ is a retract of the functor $BG\amalg e\to e\amalg e$, which is in $W^\dagger$, and is not itself in $W^\dagger$.\hfill$\square$

Consequences for the reading of the chapter. The least class satisfying Definition~1.1 is $\mathbf W_\infty$ (commentary at ms.~87), which does satisfy the relative axiom; the example shows that a larger class satisfying the definition need not. Whether Loc~(1 bis) (stability under retracts) together with (2), (3) implies (3 bis) is not decided by it. In this companion, and in the later literature, ``fundamental localizer'' without qualification means a class satisfying the relative axiom; statements proved here for fundamental localizers are not claimed for the classes of Definition~1.1.
\end{dcomment}

%%FILE I-b
%%APPEND title="Universal equivalence and the 1-complete test"
Two cautions on the repaired notions. The repair of ms.~99--100 is asymmetric: it uses constant coefficients on the total category and locally constant ones on the base only. And the fibrations defined by coefficients are not the proper functors of the localizer of locally constant coefficients: for the derivator of a nontrivial presentable coefficient theory, the class of the functors whose right Kan extension, after every base change, preserves locally constant coefficients and is computed on strict fibres, and the class of the $W^{\mathrm{lc}}_{\mathbb D}$-proper functors, are each not contained in the other (stated; the examples are not reproduced here). The inclusion of the end point $1$ in $[1]$ is proper, and its comma categories over $0\to1$ change from the empty category to a point: properness does not imply local constancy of the fibres.

%%FILE VI-c
%%BLOCK after="The final adjunction and the surviving base-change fragment"
\begin{dcomment}{status}{p204}{Base change for the two Kan extensions: the exact conditions}
The operations of this chapter can be summarized, for a derivator $\mathbb D$ on small categories in presheaf variance and a strict pullback square with $f:X\to Y$, $g:Y'\to Y$ and projections $f':X'\to Y'$, $g':X'\to X$. The comparison $f'_!g'^*\to g^*f_!$ is invertible if $f$ is smooth or $g$ is proper; the comparison $g^*f_*\to f'_*g'^*$ is invertible if $f$ is proper or $g$ is smooth (smooth and proper for $W_\infty$, as conditions on the categories of lifts). Conversely, for coefficients in spaces, $f$ is proper if the second comparison is invertible for all such squares, and dually; so these are the exact conditions. For an arbitrary pullback square neither comparison is invertible: the two end-point inclusions in $[1]$, whose pullback is empty, give a counterexample already for coefficients in finite sets, and also for stable coefficients. There is therefore no base-change formalism for all pullback squares of small categories, and the two halves of the theory (left extensions with smooth base change, right extensions with proper base change) are exchanged by passing to opposite categories: for the opposite derivator $\mathbb D^\vee(S)=\mathbb D(S^{\mathrm{op}})^{\mathrm{op}}$ one has $f_!^{\mathbb D^\vee}=((f^{\mathrm{op}})_*)^{\mathrm{op}}$ and $f_*^{\mathbb D^\vee}=((f^{\mathrm{op}})_!)^{\mathrm{op}}$. Opposition of the base exchanges the two kinds of lifting categories, hence smooth and proper; the reversal of the adjunctions needs in addition the opposite of the coefficient categories. It is a comparison of the theories over $S$ and over $S^{\mathrm{op}}$, not a duality of coefficients over a fixed base. The first two assertions were followed in outline; projection formulas and correspondences are not treated here.
\end{dcomment}
